\begin{homeworkProblem}[Seção 3-8]
  \textbf{(Unicidade da fórmula de Taylor)} Seja $f:U\subseteq\mathbb{R}^{n}\to\mathbb{R}^{m}$ ima aplicação $k$ vezes diferenciável. Sejam $\varphi_{j}\in L\left( \mathbb{R}^{n}\times\cdots\times\mathbb{R}^{n},\mathbb{R}^{m} \right)$, $j=1,2,\cdots,k$, transformações $j$-lineares e simétricas tais que
  \begin{align*}
    f(a+h)=f(a)+\varphi_{1}(h)+\varphi_2(h,h)+\cdots+\varphi_{k}(h,\cdots,h)+r_{k}(h),
  \end{align*}
  onde
  \begin{align*}
    \lim_{h \to 0}\frac{r_{k}(h)}{\norm{h}^{k}}=0.
  \end{align*}
  Então
  \begin{align*}
    \varphi_{j}=\frac{1}{j!}D^{j}f(a),\quad\text{ para todo }j=1,\cdots,k.
  \end{align*}
  \begin{solution}
    Seja $\phi:[0,1]\to\mathbb{R}^{m}$ definida por $\phi(t)=f(a+th)$  
  \end{solution}
\end{homeworkProblem}
